\chapter*{Introduction générale}
La gestion d'un entrepôt ne se limite plus à enregistrer des entrées et des sorties de marchandises. Un entrepôt moderne doit connaître l'état du stock, organiser les emplacements, préparer les commandes, coordonner les opérateurs et les équipements, prendre en compte les contraintes propres aux produits et échanger avec plusieurs systèmes de l'entreprise. Dans ce contexte, le Warehouse Management System (WMS) joue un rôle central : il structure les opérations depuis la réception jusqu'à l'expédition et assure la traçabilité des mouvements.

Un WMS traditionnel couvre déjà les fonctions essentielles : réception, contrôle, mise en stock, gestion des emplacements, inventaire, réapprovisionnement, allocation, picking, packing et expédition. La transformation vers un Smart WMS ne consiste donc pas à remplacer ce socle par une intelligence artificielle généralisée. Elle consiste plutôt à renforcer le système par des capacités activables lorsque le besoin, les données et l'infrastructure le justifient : prévision de charge, optimisation de slotting, détection d'anomalies, supervision énergétique, vision industrielle, RFID ou orchestration avec WCS/WES.

Le stage réalisé au sein de Smart Automation Technologies s'inscrit dans cette logique. L'objectif principal est de concevoir une architecture de référence capable de guider la transformation d'un WMS traditionnel vers un Smart WMS. Le travail porte sur la conception métier, fonctionnelle et applicative logique ; il ne vise pas le développement complet d'un produit WMS industriel.

\begin{problematique}
Comment concevoir une architecture Smart WMS suffisamment complète pour couvrir les fonctions essentielles d'un entrepôt, tout en restant générique, modulaire et adaptable aux contraintes produits, aux données disponibles et au niveau d'automatisation de chaque entreprise ?
\end{problematique}

La démarche adoptée suit une logique SCAN -- PLAN -- ACT. La phase SCAN a permis de comprendre les solutions WMS leaders et les limites d'un WMS classique. La phase PLAN a structuré les besoins, les use cases et l'architecture fonctionnelle cible. La phase ACT prépare la validation du modèle, la modélisation UML et une architecture applicative logique pouvant servir de base à une future implémentation.
